\documentclass[dvipdfmx, twocolumn, 10pt]{bxjsarticle}

\usepackage{geometry}
\geometry{margin=20truemm}
\usepackage{bm}
\usepackage{tabularx}
\usepackage{listings}
\lstset{
  basicstyle=\footnotesize\ttfamily,
  breaklines=true,
  breakatwhitespace=false,
  columns=fullflexible,
  keepspaces=true,
}

\title{\textbf{\LaTeX\ レポート作成ガイド}}
\author{}
\date{\today}

\begin{document}
\maketitle

\texttt{main.tex} のテンプレートをベースに, 研究報告書・論文の書き方をまとめたガイドです.

% ===========================
\section{ドキュメントクラスとレイアウト}

\begin{lstlisting}
\documentclass[dvipdfmx, twocolumn, 10pt]{bxjsarticle}
\end{lstlisting}

\begin{tabularx}{\linewidth}{lX}
\hline
オプション & 意味 \\
\hline
\texttt{dvipdfmx}    & PDF 変換ドライバ（日本語環境では必須） \\
\texttt{twocolumn}   & 2段組レイアウト \\
\texttt{10pt}        & 本文フォントサイズ \\
\texttt{bxjsarticle} & 日本語対応の article クラス \\
\hline
\end{tabularx}

\medskip
余白の調整は \texttt{geometry} パッケージで行う:

\begin{lstlisting}
\usepackage{geometry}
\geometry{margin=20truemm}
% または細かく指定
% \geometry{top=15mm, bottom=25mm, left=20mm, right=20mm}
\end{lstlisting}

% ===========================
\section{よく使うパッケージ}

\begin{lstlisting}
\usepackage{bm}    % 数式中の太字ベクトル: \bm{x}
\usepackage{tikz}  % 図の描画
\usetikzlibrary{arrows.meta, positioning, decorations.pathreplacing, calc}
\end{lstlisting}

% ===========================
\section{文書情報}

\begin{lstlisting}
\title{\textbf{タイトル}}
\author{氏名}
\date{\today}   % 今日の日付を自動挿入

\begin{document}
\maketitle      % タイトルブロックを出力
\end{lstlisting}

% ===========================
\section{要旨 (Abstract)}

\begin{lstlisting}
\begin{abstract}
研究の背景・目的・手法・結果を200〜400字程度でまとめる。
\end{abstract}
\end{lstlisting}

\textbf{書き方のポイント}
\begin{itemize}
  \item 背景 $\to$ 問題提起 $\to$ 手法 $\to$ 成果の順で書く
  \item 数式・図表への参照は入れない
  \item 専門用語を初出時に略称と併記する（例: 連合学習 (Federated Learning, FL)）
\end{itemize}

% ===========================
\section{セクション構成}

\begin{lstlisting}
\section{緒言}
\subsection{背景}
\subsubsection{詳細（必要なら）}
\end{lstlisting}

\textbf{典型的な構成}
\begin{enumerate}
  \item 緒言（背景・研究の位置づけ）
  \item 研究課題（問いの明確化）
  \item 実験計画 / 手法
  \item 評価方法
  \item 結果・考察（実験後）
  \item 結言
\end{enumerate}

% ===========================
\section{文献引用 (BibLaTeX)}

\subsection{設定}

\begin{lstlisting}
\usepackage[
  backend=biber,   % 処理エンジン（pbiber ではなく biber を使う）
  style=numeric,   % [1] 形式の番号引用
  sorting=none,    % 登場順に番号を振る
  date=year,       % 年のみ表示
  isbn=false,
  doi=true,
  giveninits,      % 名前をイニシャル表記
]{biblatex}
\addbibresource{references.bib}  % .bib ファイルを指定
\end{lstlisting}

\subsection{引用の書き方}

\begin{lstlisting}
この研究\cite{著者年キーワード}では〜が示されている.

% 複数引用
先行研究\cite{foo2020,bar2021}によると〜.
\end{lstlisting}

\subsection{文献リストの出力}

\begin{lstlisting}
\printbibliography   % 本文末尾に配置
\end{lstlisting}

\subsection{.bib ファイルの記述例}

\begin{lstlisting}
@inproceedings{mcmahan2017communication,
  author    = {McMahan, Brendan and others},
  title     = {Communication-Efficient Learning of Deep Networks},
  booktitle = {Proceedings of AISTATS},
  year      = {2017},
}

@article{li2020federated,
  author  = {Li, Tian and others},
  title   = {Federated Optimization in Heterogeneous Networks},
  journal = {Proceedings of MLSys},
  year    = {2020},
}
\end{lstlisting}

\subsection{コンパイル手順（BibLaTeX + biber）}

\begin{lstlisting}
platex main.tex
biber main
platex main.tex
platex main.tex
\end{lstlisting}

または \texttt{latexmk} を使う（推奨）:

\begin{lstlisting}
latexmk -pdfdvi main.tex
\end{lstlisting}

% ===========================
\section{図の挿入}

\subsection{1段幅の図}

\begin{lstlisting}
\begin{figure}[h]
  \centering
  \includegraphics[width=0.9\linewidth]{figures/image.png}
  \caption{図のキャプション}
  \label{fig:label}
\end{figure}
\end{lstlisting}

\subsection{2段組をまたぐ幅広の図}

\begin{lstlisting}
\begin{figure*}[t]   % * で全幅、[t] でページ上部に配置
  \centering
  % ...
  \caption{キャプション}
  \label{fig:wide}
\end{figure*}
\end{lstlisting}

\subsection{図の参照}

\begin{lstlisting}
図\ref{fig:label}に示すように〜.
\end{lstlisting}

\subsection{TikZ で図を描く（画像不要）}

\begin{lstlisting}
\begin{tikzpicture}
  \draw[-{Stealth}] (0,0) -- (5,0) node[right] {時間};
  \draw[rectangle, draw, fill=cyan!30] (1,0) rectangle (3,0.5)
    node[pos=.5] {イベント};
\end{tikzpicture}
\end{lstlisting}

% ===========================
\section{箇条書き}

\subsection{番号なし}

\begin{lstlisting}
\begin{itemize}
  \item 項目1
  \item 項目2
  \begin{itemize}         % ネストも可能
    \item サブ項目
  \end{itemize}
\end{itemize}
\end{lstlisting}

\subsection{番号あり}

\begin{lstlisting}
\begin{enumerate}
  \item 研究課題1
  \item 研究課題2
\end{enumerate}
\end{lstlisting}

% ===========================
\section{強調・テキスト修飾}

\begin{lstlisting}
\textbf{太字}          % 強調（用語の初出など）
\textit{イタリック}    % 英語の書名・変数名など
\underline{下線}
\end{lstlisting}

% ===========================
\section{数式}

\subsection{インライン数式}

\begin{lstlisting}
損失関数 $L(\theta)$ を最小化する.
\end{lstlisting}

\subsection{別行立て数式}

\begin{lstlisting}
\begin{equation}
  \hat{y} = \sigma\left(\sum_{i} w_i x_i + b\right)
  \label{eq:sigmoid}
\end{equation}

式\eqref{eq:sigmoid}より〜.
\end{lstlisting}

\subsection{太字ベクトル（bm パッケージ）}

\begin{lstlisting}
$\bm{x} \in \mathbb{R}^d$
\end{lstlisting}

% ===========================
\newpage
\section{ディレクトリ構成の例}

\begin{lstlisting}
report/
├── main.tex          # メインファイル
├── references.bib    # 参考文献データベース
├── figures/          # 画像ファイル
│   └── diagram.png
└── sections/         # セクションを分割する場合
    ├── intro.tex
    └── experiment.tex
\end{lstlisting}

分割したファイルの読み込み:

\begin{lstlisting}
\input{sections/intro}
\end{lstlisting}


% ===========================
\section{提出前チェックリスト}

\begin{itemize}
  \item 全ての \texttt{\textbackslash ref} が正しく解決されている（\texttt{??} になっていない）
  \item 全ての \texttt{\textbackslash cite} が文献リストに存在する
  \item 図・表にキャプションと \texttt{\textbackslash label} がある
  \item 要旨が研究の全体像を正確に伝えている
  \item 略称は初出時に正式名称とともに定義されている
  \item PDF に日本語フォントが埋め込まれている
  \item ページ数・フォーマットが指定要件を満たしている
\end{itemize}

\end{document}
